\section{Related Work}
\label{ch:related}

The work in this thesis sits at the intersection of three lines of research: controllable text generation, energy-based and sampling-based decoding, and the study of the geometry and behaviour of language-model representations. This section reviews each in turn. Rather than cataloguing methods, it tells the story of how the field arrived at energy-based sampling as an approach to control, because understanding that trajectory is what makes the unexamined assumption at its heart visible, and it is that assumption which the thesis tests. The section closes by stating precisely the gap in the literature that the work addresses.

\subsection{Controllable Text Generation}
\label{sec:rel-ctg}

Approaches to controllable generation can be organized by how much they change the underlying model, and reviewing them in that order reveals a progression toward ever lighter-touch intervention, of which energy-based sampling is the logical endpoint.

At one extreme are the \textbf{fine-tuning methods}, which adjust the model's parameters so that the desired behaviour is internalized. The earliest of these simply continued training on attribute-specific data, and the modern and dominant form is reinforcement learning from human feedback \citep{ouyang2022training}, which optimizes the model against a learned reward that encodes human preferences. These methods work, and the widely used aligned assistants are their fruit, but the control they produce is fixed once training is complete. A new attribute requires new data and a new training run, and the cost of that adaptation is precisely what motivates the search for methods that impose control at inference time. Conditioning methods such as CTRL \citep{keskar2019ctrl} occupy a middle ground, prepending control codes to the input so that a single model can be steered among a fixed set of attributes anticipated during training, but they too cannot accommodate a genuinely novel constraint after the fact.

At the other extreme are the \textbf{decoding-time methods}, which leave the parameters untouched and intervene only during generation, and these are the ones that promise the inference-time flexibility fine-tuning cannot offer. Plug and Play Language Models \citep{dathathri2020plug} perturb the hidden states of a frozen model at each step in the direction an attribute classifier's gradient indicates. FUDGE \citep{yang2021fudge} and GeDi \citep{krause2021gedi} instead reweight the next-token distribution using a discriminator that predicts, from a prefix, whether the eventual completion will bear the desired attribute. Effective within their scope, they share a structural limit: they remain autoregressive, modifying the local per-token decision while retaining the left-to-right order and so inheriting its inability to revise an earlier token in light of a later one. Pursuit of a method that could edit any part of a sequence in light of the whole is what leads to the sampling-based family.

\subsection{Energy-Based and Sampling-Based Decoding}
\label{sec:rel-energy}

The decisive conceptual move, and the one that defines the family of methods this thesis studies, is to abandon left-to-right generation entirely and to cast decoding as sampling from an energy defined over whole sequences. Once a sequence is treated as a single object drawn from a distribution, rather than as a stream of tokens produced in order, a global constraint can be imposed on the completed object directly, and any position in the sequence can be revised in light of any other. This is the promise that makes the approach attractive, and several influential systems realize it.

COLD decoding \citep{qin2022cold} performs energy-based constrained generation by running Langevin dynamics on a continuous relaxation of the sequence, combining a fluency energy derived from a language model with task-specific constraint energies, and demonstrated it on tasks such as constrained generation and abductive reasoning. MuCoLa \citep{kumar2022gradient} samples in the continuous embedding space and projects back to tokens; it is the closest published antecedent of the continuous sampler studied here, and its formulation together with its use of a projection to return to discrete text are exactly the ingredients whose geometric consequences Section~\ref{ch:results} examines. The continuous sampler of this thesis implements that shared mechanism faithfully, and the constrained \emph{mucola} arm of the experiments is that comparison, run at the level of the shared energy rather than of the original papers' tasks and metrics, which use different benchmarks and are out of scope. Related efforts include gradient-based inference for lexically constrained generation and infilling \citep{sha2020gradient} and discrete samplers adapted from the energy-based-model literature, among them the gradient-informed proposals of \citet{grathwohl2021oops} and the discrete Langevin proposals of \citet{zhang2022langevin} on which the discrete sampler here is directly based. A third strategy relaxes the categorical choice itself with a straight-through estimator \citep{bengio2013estimating} or the Gumbel-softmax reparameterization \citep{jang2017categorical}; it is out of scope because it alters the model's forward computation rather than leaving the frozen likelihood untouched.

One feature recurs across this literature and is rarely brought to the foreground. The Metropolis--Hastings correction, which as Section~\ref{sec:bg-mh} explained is what separates genuine sampling from biased optimization, is frequently omitted: methods are described in the language of sampling from an energy while in practice running a noise-augmented gradient descent, and where the correction is included, the difficulty of computing its reverse-proposal term on a landscape defined by projection is seldom examined. Taking the correction seriously, implementing it faithfully for both a discrete and a continuous sampler, and measuring what its inclusion reveals is one of the contributions of this thesis; what that measurement licenses one to say about published systems is treated in Section~\ref{sec:disc-related}, where the evidence for it has been presented.

Two further lines bear on the target these methods sample from. The machine-translation literature established, more sharply than anywhere else, that the mode of a well-trained sequence model is not where the quality is. \citet{stahlberg2019nmt} showed that for a large fraction of inputs the globally most probable translation under a neural model is the empty string, and \citet{eikema2020map} argued that maximum-a-posteriori decoding is therefore the wrong criterion, the model's useful information lying in the bulk of the distribution rather than at its peak; \citet{meister2023locally} give the complementary account, that high-probability continuations are dispreferred by human readers because they carry atypically little information. These are the published form of the likelihood trap measured in Section~\ref{sec:results-trap}, and they matter here because a method that succeeds at concentrating mass on the lowest-energy sequences of a frozen likelihood succeeds at reaching exactly the region these works identify as degenerate. \citet{deng2020residual} take the constructive route of learning a residual energy on top of a language model rather than treating the likelihood itself as the energy.

A distinct branch anticipates the diagnosis this thesis reaches, because not every method that samples from a frozen-language-model energy uses the \emph{gradient} of that energy. Mix-and-Match \citep{mireshghallah2022mixmatch} performs controllable generation by Metropolis--Hastings over a product-of-experts energy, proposing token substitutions from an auxiliary masked language model and scoring them with the frozen energy, which is therefore only ever evaluated. The twisted sequential Monte Carlo methods of \citet{lew2023sequential} and \citet{zhao2024probabilistic} likewise sample from an intractable posterior defined by a frozen model using importance weights and resampling rather than energy gradients. Two older members of the same branch are the direct antecedents of this study's most successful baseline: \citet{wang2019bert} sample text from a frozen masked language model by Gibbs sampling over token positions, and \citet{miao2019cgmh} perform constrained sentence generation by Metropolis--Hastings over word-level insert, delete and replace proposals. That these gradient-free samplers operate successfully on essentially the same frozen-model energies that defeat the gradient-guided methods is a strong hint about where the difficulty lies, and Section~\ref{sec:results-baselines} makes it concrete on the identical energy. Their existence is part of why the negative result of this thesis is scoped to one derivative rather than to the training-free premise.

\subsection{Amortized Sampling and GFlowNets}
\label{sec:rel-gflownet}

Generative Flow Networks were introduced by \citet{bengio2021flow} as a method for sampling discrete, compositional objects in proportion to a reward, with the explicit motivation of producing diverse high-reward samples rather than the single mode that reward-maximizing reinforcement learning tends to collapse onto. The framework was subsequently formalized and given a theoretical foundation \citep{bengio2023gflownet}, and a family of training objectives was developed to address the practical difficulty of assigning credit along long generative trajectories, including trajectory balance \citep{malkin2022trajectory} and its partial-episode sub-trajectory refinement \citep{madan2023learning}. The application most directly relevant here is that of \citet{hu2024amortizing}, who framed a range of language tasks, including infilling and reasoning, as sampling from an intractable posterior defined by a frozen language model, and trained a GFlowNet to perform that sampling, reporting that the resulting policies yield diverse samples that transfer better to downstream use than those obtained by reward maximization.

The GFlowNet component of this thesis builds directly on that formulation and codebase, using the same modified sub-trajectory objective and the same parameter-efficient adaptation of a frozen base model \citep{hu2022lora}. The question posed here is narrower: their concern is whether amortized sampling is useful, and they show that it is where the reward defines a well-behaved target, whereas the concern here is whether it \emph{escapes a specific pathology}, that of the frozen autoregressive likelihood used as an energy over free-form sequences. The two are compatible, and Section~\ref{ch:results} takes up the sharper one by direct experiment.

\subsection{The Geometry and Behaviour of Language-Model Representations}
\label{sec:rel-geometry}

A complementary line of work concerns the representations in which these samplers operate and the behaviour of the likelihood they optimize, and it supplies part of the explanation for the results. Two phenomena from this literature are directly relevant.

The first is \textbf{anisotropy}. Transformer embeddings, contextual and static alike, do not spread uniformly through the representation space but occupy a narrow cone, so two randomly chosen token embeddings have high average cosine similarity rather than being roughly orthogonal \citep{gao2019representation, ethayarajh2019contextual}. This bears on any method treating Euclidean or cosine distance in embedding space as a measure of similarity, because anisotropy compresses the dynamic range of those distances: the difference between a good and a bad token is smaller in the embedding metric than the geometry would suggest. This thesis measures the anisotropy of both models it studies, shows that it renders Euclidean distance an unreliable evaluation metric, which is one reason the experiments adopt a divergence-based measure of contextual fit instead, and adds a consequence the prior work did not draw, that the anisotropy scale differs enough between models to make a single step-size schedule non-transferable.

The second is the failure of high-probability text to be high-quality text. \citet{holtzman2020curious} showed that the most probable sequences under a neural language model are frequently degenerate, repetitive and unnatural, and that maximization-based decoding produces markedly worse text than sampling. That matters here because any procedure concentrating probability mass on the lowest-energy regions of the frozen-likelihood energy thereby concentrates it on degenerate text. This thesis reproduces the phenomenon quantitatively on its own models, terms it the \emph{likelihood trap} in the energy-based setting, and uses it to argue that the objective the samplers pursue is itself undesirable, quite apart from whether they can pursue it effectively.

\subsection{Diffusion Models for Text and Bidirectional Denoising}
\label{sec:rel-diffusion}

A final body of work concerns a comparative model family rather than the methods this thesis evaluates. Diffusion models, which generate by gradually reversing a noising process, have become the dominant approach to image generation and have been adapted to text in two broad families. The first embeds tokens into a continuous space and runs a continuous diffusion there, as in Diffusion-LM \citep{li2022diffusion}, which demonstrated controllable generation by applying classifier guidance to the continuous latents at each denoising step. The second operates directly on discrete tokens, as in the score-entropy discrete diffusion of \citet{lou2024discrete}, which learns the ratios of the data distribution over a discrete state space and has narrowed much of the quality gap between diffusion and autoregressive language models.

Two things distinguish these models from an autoregressive one, and keeping them apart matters for how the comparison in Section~\ref{sec:results-diffusion} may be read. The first is the objective: by the equivalence of \citet{vincent2011connection}, learning to denoise is learning the score of the data distribution, which is what gives score-based generative modeling its name and its mechanism \citep{song2019generative, ho2020denoising}. The second is the conditioning direction: a denoising step reads the surviving context on both sides of the position it fills, which is the query an in-place revision poses and which no left-to-right factorization affords. A diffusion model therefore differs from an autoregressive one along both axes at once, alongside differences in scale and corpus, so it is a comparative family and not by itself a control that isolates either axis. A bidirectional masked language model, which has the second property without the first, is what separates them, and Section~\ref{sec:results-mlm} runs it.

\subsection{The Gap Addressed by This Thesis}
\label{sec:rel-gap}

The literature above establishes energy-based sampling as an attractive and actively pursued route to controllable generation, and supplies both the samplers this thesis studies and the amortized alternative it examines. What it does not supply is a direct, controlled test of the assumption the enterprise rests on, that the gradient of a frozen autoregressive likelihood is a usable search direction on the discrete text manifold, together with a mechanistic account of what happens when that assumption fails. Negative observations, where they appear, are reported in passing as brittleness, hyperparameter sensitivity, or a need for careful tuning, rather than isolated and explained. This thesis addresses that gap. It designs an ablation that separates the direction of the gradient from its magnitude and so tests the assumption rather than a proxy for it; it constructs diagnostics that identify the geometric, statistical and analytic mechanisms behind the result; and it locates the deficiency precisely, in the coordinates the derivative is taken in and in what the resulting proposal may condition on, rather than in the sampler, the energy, or the objective the model was trained with.
